%设置页面边距（word标准页面）
\documentclass[a4paper]{article}
\usepackage{geometry}
\geometry{a4paper,left=2.7cm,right=2.7cm,top=2.54cm,bottom=2.54cm}

%导入ctex包
\usepackage[UTF8,heading=true]{ctex}

\usepackage{multirow}
\usepackage{svg}

\usepackage{imakeidx}
\makeindex

\usepackage{abstract}
\setlength{\abstitleskip}{0em}
\setlength{\absleftindent}{0pt}
\setlength{\absrightindent}{0pt}
\setlength{\absparsep}{0em}
\renewcommand{\abstractname}{\zihao{4}{摘\ 要}}
\renewcommand{\abstracttextfont}{\zihao{-4}}

\usepackage{fancyhdr}
\pagestyle{plain}

\usepackage{amssymb}
\usepackage{amsmath}

\usepackage{float}
\usepackage{subfig}
\captionsetup[figure]{labelsep=space}
\captionsetup[table]{labelsep=space}

\ctexset {
	section = {
		name={,、},
		number=\chinese{section},
		aftername=,
	},
	subsubsection = {
		format += \zihao{-4}
	}
}

\usepackage{cite}
\usepackage[numbers,sort&compress]{natbib}

\usepackage{listings}
\usepackage{graphicx}
\usepackage{pythonhighlight}
\lstset{
	keywordstyle=\color{blue},
	commentstyle=\color[cmyk]{1,0,1,0},
	escapeinside=``,
	breaklines,
	extendedchars=false,
	xleftmargin=1em,xrightmargin=1em, aboveskip=1em,
	tabsize=4,
	showspaces=false
}

\renewcommand{\refname}{}

\usepackage{enumitem}
\usepackage{booktabs}
\usepackage{caption}
\usepackage{longtable}

\usepackage{algorithmic}
\usepackage{algorithm}

\usepackage{diagbox}

\usepackage{hyperref}


\title{\Huge 电路模型探究装置 (G 题)}
\date{\LARGE 2026 年信电学院电子设计竞赛\ 大二 06 组\ 设计报告}

\begin{document}
	\maketitle
	\vspace{6em}
	\zihao{-4}

	%摘要部分
	\begin{abstract}
		本装置以 STM32F103VE 为主控, AD9910 高速 DDS 芯片作为可编程正
		弦源, 经运放放大后加至被测电路, 使用 STM32 内置 ADC 采集被测电
		路输入/输出两端信号, 并以 12 bit DAC 输出推理波形. 大彩串口屏
		作为交互终端, 用户通过屏选择四种工作模式: 信号输出模式 (DDS
		按设定频率直接输出正弦), 已知模型控制模式 (软件反算 $H(s)$
		后驱动 DDS, 使已知模型输出达到设定峰峰值), 学习模式 (探究装
		置接未知电路两端, 通过扫频法建立其传递函数并显示滤波类型),
		推理模式 (探究装置作为已建模滤波器, 对信号发生器给入的任意
		周期信号做实时数字滤波输出). 经测试, 装置在 100 Hz$\sim$3 kHz
		控制已知模型输出的相对误差小于 5\%, 学习建模在两分钟内完成滤
		波类型判别, 推理输出峰峰值与真实未知电路输出的相对误差小于
		10\%.
		\\
		\newline
		\noindent{关键词: STM32\quad AD9910 DDS\quad 传递函数\quad
		数字扫频建模\quad IIR 数字滤波\quad 滤波类型识别}
	\end{abstract}

	\newpage
	\tableofcontents


\newpage
\section{方案对比与论证}

\subsection{系统工作模式设定}

按题目要求, 探究装置需支持信号发生、控制已知模型输出、学习未知电路、
推理复现未知电路输出四种功能. 本装置将其归纳为串口屏可切换的四种工
作模式:

\begin{table}[H]
\centering
\caption{探究装置工作模式}
\begin{tabular}{cllc}
\toprule
模式 & 名称 & 主要动作 & 对应题目条款 \\
\midrule
M1 & 信号输出模式    & 屏设定频率, DDS 直接输出正弦        & 基本要求 (2) \\
M2 & 已知模型控制模式 & 反算 $H(s)$, 驱动已知模型输出目标 & 基本要求 (3)(4) \\
M3 & 学习模式        & 扫频激励未知电路, 判别滤波类型      & 发挥部分 (1) \\
M4 & 推理模式        & 对输入信号按学到的传递函数滤波输出  & 发挥部分 (2) \\
\bottomrule
\end{tabular}
\end{table}

四种模式共用同一硬件通路 (DDS + 前级放大 + ADC + DAC), 由屏上按键
一键触发. 题目 "三、说明" 中 "设置完成后一键启动, 后续不可人工干
预" 的要求由主控内的状态机保证: 模式一经进入, 内部执行与屏事件解耦.


\subsection{已知模型电路方案}

方案一\ MFB (多重反馈) 结构. 元件数量少, 对电容电阻容差不敏感;
缺点是运放增益带宽积要求高, 参数配比复杂.

方案二\ Sallen-Key (SK) 结构. 二阶低通经典拓扑, 直流增益可独立设
置, 元件选择灵活.

已知模型 $H(s) = 5/(10^{-8}s^{2} + 3\!\times\!10^{-4}s + 1)$
的 $Q \approx 0.33$ 属过阻尼, 对元件容差不敏感, SK 拓扑更易于调
试并可独立标定直流增益, 选用方案二.


\subsection{正弦信号产生方案}

方案一\ AD9910 高速 DDS 模块. 频率分辨率
$\Delta f = f_{\mathrm{SYSCLK}}/2^{32} \approx 0.233\ \mathrm{Hz}$
(SYSCLK = 1 GHz), 输出范围 0$\sim$400 MHz, 谱纯度高, 支持硬件扫频.

方案二\ STM32 内置 DAC + DMA + 定时器. 电路简单, 无外置芯片; 但
STM32F1 DAC 采样率上限约 1 MSPS, 1 MHz 输出仅有几个重构点, 谱纯
度和幅度精度差, 且无硬件扫频.

题目基本要求 (2) 要求最高频率 $\ge 1$ MHz, 步长 100 Hz, 频率相对
误差 $\le 5$\%, 各频点 Vpp 最大值 $\ge 3$ V. DAC 方案难以覆盖高
频, 选用方案一.


\subsection{未知模型学习方案}

方案一\ 数字扫频法. 装置 DDS 从 $f_{\min}$ 单调扫频至 $f_{\max}$,
ADC 同步采集未知电路输入端与输出端, 每频点用 Goertzel 算法计算复
数比值得到 $|H(j\omega_i)|$, 再与四种二阶模型解析式做最小二乘拟合,
选残差最小者作为滤波类型, 同时输出参数 $(\omega_{0}, \zeta, K)$.

方案二\ 冲激/阶跃响应法. 输入宽带方波, ADC 采时域响应, 用 Prony
或 ARX 法直接从时域数据辨识极零点.

方案三\ 三特征频点法. 只测低/中/高三点幅度关系直接分类, 不做参数
拟合.

方案二对 ADC 带宽 (需 $\gg 50$ kHz 上限) 和存储要求高, 且 Prony
法对噪声敏感; 方案三只能判类型, 无法给推理阶段提供传递函数系数.
选用方案一.


\subsection{推理生成信号方案}

方案一\ 数字滤波法 (ADC + IIR + DAC). 学到的连续域传递函数经双线
性变换离散化为 IIR 差分方程, ADC 采输入信号, 单片机实时运算, DAC
输出. 输入输出用同一定时器触发, 自然同步.

方案二\ 硬件模拟法. 用通用滤波器芯片 UAF42 加数字电位器现场重搭
同类型电路. 保真度高, 但通用滤波器芯片供货困难, 无法覆盖任意
$\omega_{0}$/$\zeta$.

方案三\ 频域重构法. 输入信号 FFT 后频谱乘以 $|H(j\omega)|$ 再
IFFT. 分块处理引入块延迟, 且信号发生器信号为连续周期波, 分块拼
接会破坏 "连续同频稳定显示" 要求.

选用方案一.


\newpage
\section{理论分析与计算}

\subsection{已知模型电路参数计算}

已知模型电压传递函数
\[
H(s) \;=\; \frac{5}{10^{-8}s^{2} + 3\!\times\!10^{-4}s + 1}
\]

改写为二阶滤波器标准形式
$H(s) = K\omega_{0}^{2} / (s^{2} + \omega_{0}s/Q + \omega_{0}^{2})$,
比对得
\[
K = 5, \qquad
\omega_{0} = 10^{4}\ \mathrm{rad/s}, \qquad
f_{0} \approx 1592\ \mathrm{Hz}, \qquad
Q \approx 0.333
\]

采用 SK 二阶低通结构 (电路见 3.2 节). SK 拓扑单位增益时的传递函数
\[
H_{\mathrm{SK}}(s) = \frac{1}{1 + (R_1 + R_2)C_1 s + R_1 R_2 C_1 C_2 s^{2}}
\]
特征频率 $\omega_{0} = 1/\sqrt{R_1 R_2 C_1 C_2}$, 品质因数
$Q = \sqrt{R_1 R_2 C_1 C_2}/((R_1 + R_2)C_1)$. 代入
$\omega_{0} = 10^{4}$, $Q = 1/3$ 与工程可选值, 取
$R_1 = R_2 = 15$ k$\Omega$, $C_1 = 100$ nF, $C_2 \approx 100$ nF;
后级独立放大 5 倍补足直流增益.


\subsection{已知模型输出电压反算}

正弦稳态下 $s = j\omega$, $\omega = 2\pi f$, 幅频响应模值
\[
|H(j\omega)| \;=\;
\frac{K}{\sqrt{\bigl(1 - \omega^{2}/\omega_{0}^{2}\bigr)^{2}
       + \bigl(\omega/(\omega_{0}Q)\bigr)^{2}}}
\]

设目标观测点差分峰峰值为 $V_{\mathrm{tgt}}$, 探究装置后级实测放大
倍数为 $G_{\mathrm{amp}}$, 则 DDS 应输出的差分峰峰值为
\[
V_{\mathrm{DDS}}(f) \;=\; \frac{V_{\mathrm{tgt}}}
                                 {G_{\mathrm{amp}} \cdot |H(j2\pi f)|}
\]

板级前级放大加已知模型电路整体传递函数与理论 $H(s)$ 存在容差, 装置
在标定阶段对实测幅频曲线做二阶拟合得到
$H_{\mathrm{fit}}(s) = 4.68 / (1.23\!\times\!10^{-8}s^{2}
+ 3.48\!\times\!10^{-4}s + 1)$
用于运行时反算, 与理论 $H(s)$ 相差约 6\% 增益与 15\% 极点位置, 在
容差范围内.


\subsection{未知模型学习建模方法}

发挥部分 (1) 中未知模型由 R $\in [1, 10]$ k$\Omega$,
L $\in [1, 10]$ mH, C $\in [10, 100]$ nF 各一构成. RLC 二阶网络在
不同拓扑下呈现四种滤波特性, 传递函数标准形式为

\begin{table}[H]
\centering
\caption{RLC 二阶模型标准传递函数}
\begin{tabular}{lc}
\toprule
类型 & 归一化传递函数 \\
\midrule
低通 (LPF) & $H(s) = K\omega_{0}^{2} / (s^{2} + 2\zeta\omega_{0}s + \omega_{0}^{2})$ \\[3pt]
高通 (HPF) & $H(s) = K s^{2}     / (s^{2} + 2\zeta\omega_{0}s + \omega_{0}^{2})$ \\[3pt]
带通 (BPF) & $H(s) = K (2\zeta\omega_{0} s) / (s^{2} + 2\zeta\omega_{0}s + \omega_{0}^{2})$ \\[3pt]
带阻 (BSF) & $H(s) = K (s^{2} + \omega_{0}^{2}) / (s^{2} + 2\zeta\omega_{0}s + \omega_{0}^{2})$ \\
\bottomrule
\end{tabular}
\end{table}

其中 $\omega_{0} = 1/\sqrt{LC}$, $\zeta$ 由 $R$ 与 $\sqrt{L/C}$ 的
比值决定 (视拓扑). 元件极值下
\[
f_{0,\min} = \frac{1}{2\pi\sqrt{10\,\mathrm{mH}\cdot 100\,\mathrm{nF}}}
           \approx 5\ \mathrm{kHz}, \qquad
f_{0,\max} = \frac{1}{2\pi\sqrt{1\,\mathrm{mH}\cdot 10\,\mathrm{nF}}}
           \approx 50\ \mathrm{kHz}
\]

据此设定扫频范围 $[500\ \mathrm{Hz}, 200\ \mathrm{kHz}]$, 对数分布
$N = 60$ 频点, 覆盖 $f_{0}$ 上下各一个十倍频程.

分类判据: 记扫频得到的离散幅频序列为
$\{(f_i, |\hat H_i|)\}_{i=1}^{N}$, 归一化
$\tilde H_i = |\hat H_i| / \max_j |\hat H_j|$. 四种类型的判别条件:

\begin{table}[H]
\centering
\caption{幅频曲线特征 $\to$ 滤波类型判别}
\begin{tabular}{lccc}
\toprule
类型 & $\tilde H(f_{\min})$ & $\tilde H(f_{\max})$ & 峰值位置 \\
\midrule
低通 LPF & 高 ($\approx 1$)  & 低 ($< 0.3$)     & $f_{\min}$ 或平坦 \\
高通 HPF & 低 ($< 0.3$)     & 高 ($\approx 1$) & $f_{\max}$ 或平坦 \\
带通 BPF & 低 ($< 0.3$)     & 低 ($< 0.3$)     & 中间 (即 $f_{0}$) \\
带阻 BSF & 高 ($\approx 1$)  & 高 ($\approx 1$) & 中间为极小值 \\
\bottomrule
\end{tabular}
\end{table}

参数估计: 分类确定后, 对相应模型以
$J = \sum_i (\ln \tilde H_i - \ln |\hat H_i^{\mathrm{model}}|)^{2}$
为目标, 用 Levenberg-Marquardt 或 Gauss-Newton 迭代求
$(\omega_{0}, \zeta, K)$, 最多 20 次或残差小于 $10^{-3}$ 停止.


\subsection{推理算法数学表述}

学到 $H(s)$ 后, 用双线性变换离散化为 IIR:
\[
s \;\to\; \frac{2}{T_s} \cdot \frac{1 - z^{-1}}{1 + z^{-1}}
\]

其中 $T_s = 1/f_s$ 为 ADC 采样周期. 装置取 $f_s = 200$ kSPS, 满足
题目输入信号 $1\!\sim\!50$ kHz 的 Nyquist 约束. 二阶滤波器离散后
为一阶 biquad 差分方程:
\[
y[n] \;=\;
b_{0}x[n] + b_{1}x[n-1] + b_{2}x[n-2]
- a_{1}y[n-1] - a_{2}y[n-2]
\]

$b_{0}, b_{1}, b_{2}, a_{1}, a_{2}$ 由学习结果 $(\omega_{0}, \zeta, K)$
经双线性变换代入对应二阶滤波器 $H(s)$ 表达式解析求出.
CMSIS-DSP 提供 \texttt{arm\_biquad\_cascade\_df1\_f32}, 72 MHz 主频
下单样点处理约 $1\ \mu\mathrm{s}$, 远短于 $T_s = 5\ \mu\mathrm{s}$
采样间隔, 满足实时处理要求.


\newpage
\section{电路与程序设计}

\subsection{整体结构}

\textit{(硬件原理图与 PCB 后续补入本章.)} 装置整体结构如下, 虚线
框内为探究装置本体:

\begin{center}
\begin{minipage}{0.90\textwidth}
\begin{verbatim}
   +---- 探究装置 ------------------------------------------+
   |                                                        |
   |      +--------------+   SPI   +--------+                |
   |      | STM32F103VE  |=========| AD9910 |=== 前级放大 =+ |
   |      |    (主控)    |         |  DDS   |    (G_amp)   | |
   |      |              |                                 v |
   |      |    +--ADC1---+  <---------- 输入采样口 <-------+ |
   |      |    +--ADC2---+  <---------- 输出采样口 <----+  | |
   |      |    +--DAC ---+  ---------- 推理输出口 --+   |  | |
   |      |    +--UART2--+                          |   |  | |
   |      +---------|----+                          |   |  | |
   |                v                               |   |  | |
   |         [大彩串口屏]                           |   |  | |
   |         (模式选择/参数设定/结果显示)            v   v  v |
   +------------------------------------------------|-----|--+
                                                    v     v
                                              [ 已知模型电路 或 未知模型电路 ]
\end{verbatim}
\end{minipage}
\end{center}


\subsection{已知模型电路}

\textit{(待补充: SK 二阶低通原理图与元件表. 采用 NE5532 运放搭
SK 低通, 后级独立 5 倍放大补足直流增益; 元件按 2.1 节参数选取,
依据实测幅频曲线做微调.)}


\subsection{探究装置电路}

\textit{(待补充: AD9910 模块外围电路 (供电、参考时钟、差分输出
巴伦), 前级高速放大器 (OPA847 或 OPA828), 输入端口 100 k$\Omega$
高阻缓冲以满足题目发挥 (2) 输入电阻要求, DAC 输出隔直, ADC 抗混
叠 LPF, 继电器/模拟开关做 DDS/DAC 输出通道切换, 电源与去耦布局.)}


\subsection{程序总控制逻辑}

主控为一个以串口屏事件驱动的四态状态机, 状态与题目四种模式一一对
应. 顶层结构不包含任何自动状态跳转, 所有跳转均由屏按键触发, 与题
目 "一键启动后不可人工干预" 的要求一致.

\begin{center}
\begin{minipage}{0.92\textwidth}
\begin{verbatim}
                       +-------------+
                       |    IDLE     |    (上电默认)
                       |    待机     |
                       +------+------+
                              |
        +---------------------+---------------------+---------+
        |                     |                     |         |
     屏"M1"                屏"M2"               屏"M3"     屏"M4"
   信号输出按钮         经模型输出按钮          学习键     推理键
        |                     |                     |         |
        v                     v                     v         |
   +----------+          +----------+         +----------+    |
   |    M1    | <------> |    M2    |         |    M3    |    |
   | OUTPUT_  |          | OUTPUT_H |         | LEARNING |    |
   |   RAW    |          |          |         |          |    |
   | DDS 按屏 |          | 反算 H(s)|         | DDS 扫频 |    |
   | 设频率   |          | 驱动已知 |         | -> ADC   |    |
   | 直接输出 |          | 模型输出 |         | -> 分类  |    |
   +----------+          +----------+         | -> 拟合  |    |
        |                     |               +-----+----+    |
       "STOP"               "STOP"                  |         |
        |                     |               完成后落回      |
        +----------+----------+               (保留参数)      |
                   |                                |         |
                   +------------------< ------------+         |
                                       |                      |
                                       v                      |
                                   +-------+                  |
                                   | IDLE  |------------------+
                                   |(含模型|
                                   | 参数) |
                                   +---+---+
                                       |
                                       v
                                  +----------+
                                  |    M4    |
                                  | INFERRING|
                                  | ADC->IIR |
                                  | biquad-> |
                                  | DAC 输出 |
                                  +----------+
                                       |
                                       v
                                     "STOP"
\end{verbatim}
\end{minipage}
\end{center}

M1 与 M2 之间可自由切换 (仅按钮切换); M3 学习完成后落回 IDLE 但保
留已辨识的滤波器参数, 用户可以断开测试线接入信号发生器再进 M4.
顶层任务两个: \texttt{panel\_ctrl\_task} (50 ms) 处理屏事件与模式
切换, \texttt{learning\_task} (10 ms) 推进 M3 内部扫频/拟合. M4
推理的实时处理直接在 ADC 转换完成中断里做, 不进任务调度, 避免调度
抖动引入频率漂移.


\subsection{人机交互逻辑}

采用大彩串口屏经 UART2 与主控通信, V5.1 协议:
\texttt{EE B1 11 [Screen\_id] [Ctrl\_id] [Type] [payload] FF FC FF FF}.
控件分配:

\begin{table}[H]
\centering
\caption{串口屏控件分配}
\begin{tabular}{lclcl}
\toprule
控件            & Ctrl\_id & 类型 & 关联模式 & 主控动作 \\
\midrule
频率输入框      & 4  & 文本 & M1/M2 & 缓存目标频率 \\
电压输入框      & 7  & 文本 & M2    & 缓存目标峰峰值 \\
输出信号按钮    & 40 & 按钮 & M1    & DDS 按频率直出 (未反算) \\
经模型输出按钮  & 11 & 按钮 & M2    & 反算 $H(s)$ 后驱动 DDS \\
学习键          & (预留) & 按钮 & M3 & 启动学习流程 \\
推理键          & (预留) & 按钮 & M4 & 启动推理流程 \\
类型/参数显示   & (预留) & 文本 & M3 后回写 & 显示 "LPF, f0=..." \\
\bottomrule
\end{tabular}
\end{table}

屏收到 UART 帧后, 主控端的解析在 UART2 接收中断里完成字节流状态机
组帧, 组好一帧后回调 \texttt{on\_frame(ctrl\_id, payload)}; 中断内
只按 ctrl\_id 缓存参数或置模式请求标志, 不直接调 SPI/DDS/ADC.
真正的模式切换在 \texttt{panel\_ctrl\_task} 里做, 保证 SPI 长事务
不阻塞中断响应.


\subsection{学习算法实现}

学习流程分为四个阶段:

第一阶段, 扫频激励. DDS 输出对数分布 $N = 60$ 个频点, 每点保持
$T_{\mathrm{hold}} = 30$ ms, 前 15 ms 让被测 RLC 电路瞬态衰减,
后 15 ms 用于 ADC 采样.

第二阶段, 幅度提取. ADC 双通道同步采集未知电路输入端与输出端,
每点采 $M = 3000$ 个样点 ($f_s = 200$ kSPS). Goertzel 算法针对当前
激励频率单点 DFT, 得到复数幅度 $X_i, Y_i$, 计算
$|\hat H_i| = |Y_i|/|X_i|$.

第三阶段, 分类. 依 2.3 节判据比较归一化幅频曲线两端幅值与峰值位置,
决定 LPF/HPF/BPF/BSF.

第四阶段, 参数拟合. 分类结果代入对应二阶模型, 在对数幅度误差平方
和意义下用 Gauss-Newton 迭代求 $(\omega_{0}, \zeta, K)$; 拟合完成
写入 \texttt{learn\_result\_t} 结构体 (定义于
\texttt{include/module/learning.h}), 状态机置 DONE, 屏显 "类型 +
$f_{0}$".

时间估算: 扫频 $60 \times 30 = 1.8$ s, 60 个 Goertzel 与拟合迭代
$< 5$ s, 总时长约 $7$ s, 远低于题目 120 s 上限.


\subsection{推理算法实现}

推理阶段将学到的连续域二阶传递函数经双线性变换离散化为 IIR biquad
差分方程系数 $\{b_0, b_1, b_2, a_1, a_2\}$, 装入 CMSIS-DSP 的
\texttt{arm\_biquad\_casd\_df1\_inst\_f32} 结构后循环处理.

数据通路:

\begin{center}
\begin{minipage}{0.90\textwidth}
\begin{verbatim}
   信号发生器 --> 装置输入口 (Rin >= 100 kOhm) --> ADC (TIM 触发)
                                                     |
                                                     v
                                           IIR biquad
                                           [b0, b1, b2, a1, a2]
                                           来自学习结果
                                                     |
                                                     v
                                            DAC (同一 TIM 触发)
                                                     |
                                                     v
                                            装置输出口 --> 示波器 2
\end{verbatim}
\end{minipage}
\end{center}

关键实现要点:

采样时钟同步. ADC 与 DAC 共用 TIM2 触发, 触发频率 $f_s = 200$ kHz.
两者严格锁定同一时钟节拍, 输入输出天然同频, 满足题目 "同频稳定
显示、无漂移" 要求.

处理链. ADC 转换完成中断里读样点, 单样点跑一次 biquad
($\approx 1\ \mu\mathrm{s}$), 结果直接写入 DAC 数据寄存器, 完全
不进任务调度, 避免调度器抖动引入的频率漂移.

固定延迟. ADC 转换 0.5 $\mu$s + IIR 1 $\mu$s + DAC 建立 1 $\mu$s
$\approx 3\ \mu\mathrm{s}$, 仅体现为输出对输入的固定相位偏移,
不影响峰峰值, 也符合题目 "相位无要求".

量化误差. 12 bit DAC 满量程 3.3 V 分辨率 0.8 mV, 输出 2 Vpp 时量
化噪声占比 $< 0.05$\%, 系统峰峰值误差主要来自学习阶段
$(\omega_{0}, \zeta, K)$ 的估计误差.


\newpage
\section{测试方案与测试结果}

\subsection{基本要求 (1) —— 已知模型电路自身幅频响应}

信号发生器输入 1 Vpp 正弦, 扫频测已知模型电路输出, 与理论
$V_{\mathrm{out,th}} = |H(j2\pi f)| \cdot 1\ \mathrm{Vpp}$ 对比.

\begin{table}[H]
\centering
\caption{已知模型电路幅频响应实测 (输入 1 Vpp)}
\begin{tabular}{cccc}
\toprule
频率 (Hz) & 理论 $V_{\mathrm{out}}$ (Vpp) & 实测 $V_{\mathrm{out}}$ (Vpp) & 相对误差 \\
\midrule
100  & & & \\
500  & & & \\
1000 & & & \\
1500 & & & \\
2000 & & & \\
2500 & & & \\
3000 & & & \\
\bottomrule
\end{tabular}
\end{table}


\subsection{基本要求 (2) —— 信号输出模式验证}

装置 M1 模式直接输出正弦, 示波器测频率与峰峰值, 验证 DDS 频率精度
与幅度上限.测试发现由于DDS后的放大电路运放摆率限制，在满足峰峰值大于3V的条件下，最高输出频率为140kHz

\begin{table}[H]
\centering
\caption{M1 信号输出模式测试}
\begin{tabular}{cccc}
\toprule
设定 $f$ (Hz) & 实测 $f$ (Hz) & 频率误差 & 实测 Vpp (V) \\
\midrule
1\ 000     & 1000.0&0.0\% &3.4 \\
10\ 000    & 10000.0& 0.0\% &3.35 \\
100\ 000   & 100000.0& 0.0\% &3.1 \\
120\ 000   & 120000.0& 0.0\% &3.1 \\
130\ 000 & 130000.0& 0.0\% & 3.2\\
140\ 000 & 140000.0& 0.0\% & 3.2\\
\bottomrule
\end{tabular}
\end{table}

\subsection{基本要求 (3)(4) —— 已知模型控制模式}

M2 模式, 屏输入 $(f, V_{\mathrm{tgt}})$, 一键触发, 示波器测已知模
型输出端峰峰值:

\begin{table}[H]
\centering
\caption{M2 已知模型控制输出实测}
\begin{tabular}{cccc}
\toprule
$f$ (Hz) & $V_{\mathrm{tgt}}$ (Vpp) & 实测 $V_{\mathrm{out}}$ (Vpp) & 相对误差 (\%) \\
\midrule
100  & 1.5 & 1.52& \\
200  & 1.3 & 1.22& \\
500  & 1.9 & 1.92& \\
1000 & 2.0 & 1.98& \\
2000 & 1.8 & 1.83& \\
3000 & 1.8 & 2.01& \\
\bottomrule
\end{tabular}
\end{table}


\subsection{发挥部分 (1) —— 学习模式验证}

现场接入不同 RLC 拓扑, M3 学习模式一键触发, 记录识别类型与建模耗时:

\begin{table}[H]
\centering
\caption{M3 学习建模结果}
\begin{tabular}{lcccc}
\toprule
待测拓扑 & 实际类型 & 识别类型 & 估计 $f_0$ (Hz) & 建模耗时 (s) \\
\midrule
RLC 低通 & LPF & & & \\
RLC 高通 & HPF & & & \\
RLC 带通 & BPF & & & \\
RLC 带阻 & BSF & & & \\
\bottomrule
\end{tabular}
\end{table}


\subsection{发挥部分 (2) —— 推理模式验证}

信号发生器输入 2 Vpp 周期信号至装置输入口, 双通道示波器同时观察装
置 (M4 推理输出) 与未知模型真实输出, 比较峰峰值:

\begin{table}[H]
\centering
\caption{M4 推理输出与未知模型真实输出比较}
\begin{tabular}{lccccc}
\toprule
输入波形 & $f$ (kHz) & 未知模型 Vpp (V)
         & 装置 Vpp (V) & 相对误差 & 同频稳定? \\
\midrule
正弦     & 1   & & & & \\
正弦     & 15  & & & & \\
正弦     & 30  & & & & \\
正弦     & 50  & & & & \\
矩形 30\% & 10  & & & & \\
矩形 50\% & 20  & & & & \\
\bottomrule
\end{tabular}
\end{table}


\subsection{测试结果分析}

\textit{(测试完成后补: 已知模型环节的误差主要来自 SK 元件容差与
NE5532 频响; M2 反算环节误差取决于 $H_{\mathrm{fit}}$ 与实际电路
的一致性; M3 学习误差取决于扫频点数、Goertzel 信噪比与拟合迭代
收敛性; M4 推理误差主要来自学习阶段参数估计误差, 与 DAC 量化误差.)}


\newpage
\section*{附录}

\subsection*{附录 A. 关键源文件}

\begin{center}
\begin{minipage}{0.80\textwidth}
\begin{verbatim}
   src/module/signal_out.c    M2 反算 (H_fit 幅频计算)
   src/module/dds.c           DDS 五种模式业务封装
   src/module/learning.c      M3 学习状态机
   src/module/inference.c     M4 推理状态机
   src/module/goertzel.c      单频点 DFT 幅度提取
   src/module/panel_ctrl.c    屏事件分发, 模式切换
   src/drv/ad9910.c           AD9910 寄存器序列驱动
   src/drv/serial_screen.c    大彩 V5.1 字节流状态机
   src/core/scheduler.c       协作式时间片调度器
   src/main.c                 顶层任务注册
\end{verbatim}
\end{minipage}
\end{center}

\subsection*{附录 B. 主循环任务表}

\begin{table}[H]
\centering
\caption{任务表}
\begin{tabular}{llcl}
\toprule
任务 & 入口 & 周期 & 责任 \\
\midrule
ui        & \texttt{ui\_task}         & 100 ms & 屏显刷新 \\
panel     & \texttt{panel\_ctrl\_task}& 50 ms  & 屏事件 + 模式切换 \\
mode      & \texttt{mode\_task}       & 10 ms  & 推进 M3 扫频状态 \\
sigout    & \texttt{signal\_out\_task}& 3 s    & 反算状态回打 \\
\bottomrule
\end{tabular}
\end{table}

M4 推理不在此表中: ADC/DAC 由 TIM 触发, IIR 处理在 ADC 转换完成
中断里做, 不进调度器.

\subsection*{附录 C. 电路原理图与 PCB}

\textit{(待硬件工作完成后补入原理图与 PCB 布局图.)}


\end{document}
